% Chapter 3

\chapter{系统设计与实现}

\section{任务分析与总体架构设计}

\subsection{研究目标与输入输出定义}

本研究的核心目标是构建一套面向特定艺术风格的视频生成实验流程，并在有限算力条件下验证基于LoRA的参数高效微调方法在风格学习任务中的可行性。结合当前项目实施情况，可以将整个研究任务划分为两个层次：第一层是图像风格学习与验证，即利用特定风格图像数据对Stable Diffusion底模进行LoRA微调，并通过推理结果、对比图和可视化指标分析风格注入效果；第二层是在图像风格学习结果的基础上，进一步扩展到短视频生成原型验证，通过逐帧递推生成和时序一致性分析，探索该风格适配器在视频场景下的延展能力。

从系统输入输出角度看，本研究的输入主要包括三部分：一是原始风格图像数据及其对应文本描述，二是视频阶段所需的帧序列或文本提示词，三是云端实验配置参数，如分辨率、训练步数、LoRA秩、推理步数和视频帧数等。系统输出则包括训练得到的LoRA权重文件、单图推理结果、baseline与LoRA的对比图、批量视频片段、时序一致性评估结果以及实验过程中的配置与统计清单。通过对输入与输出的明确界定，可以把研究的实现流程组织为一个可复现、可扩展、可比较的实验链路。

\subsection{系统整体流程设计}

围绕特定风格视频生成任务，本文将系统流程划分为数据准备、模型训练、结果生成和实验评估四个核心环节。首先，在数据准备阶段，对原始风格图像进行扫描、完整性检查、中心裁剪、尺寸统一、训练验证集划分以及文本描述整理；对于视频阶段，则对连续帧序列进行抽样、裁剪缩放和片段级元数据组织。其次，在模型训练阶段，以Stable Diffusion系列模型为预训练底座，在U-Net注意力层中注入LoRA适配器，通过最小化噪声预测误差完成风格学习。随后，在结果生成阶段，分别执行图像推理、基线模型与LoRA模型的横向对比，以及短视频原型生成。最后，在实验评估阶段，结合定性观察、图像差异可视化和无参考时序一致性指标，对模型性能进行综合分析。

从实现层次上看，本文采用先图像、后视频的渐进式路线。图像阶段重点解决风格数据的规范化组织、LoRA轻量化微调以及风格注入效果验证等问题；视频阶段则在既有图像风格学习结果的基础上，进一步探索连续帧场景下的风格保持与时序稳定性\cite{ref35}。考虑到视频扩散模型训练对算力和时间成本要求较高，本文当前采用基于图像扩散模型的短视频原型方案，通过首帧生成与递推式img2img组合构造短视频片段，对连续生成的可行性进行验证。

\section{数据集构建与预处理模块设计}

\subsection{风格图像数据集构建}

数据集构建是风格学习任务中最基础也最关键的环节之一。针对本研究特定风格视频生成模型研究的目标，图像阶段首先选择具有明确风格一致性的动画风格图像作为训练数据来源，并采用与样本同名的文本文件保存caption描述。这样做的主要考虑在于：一方面，统一风格的数据能够帮助LoRA模型更加集中地学习目标画风中的配色、光影、场景构图和笔触特征；另一方面，成对的图文描述可以保证模型在训练过程中同时建立风格特征与文本条件之间的联系，提高后续提示词驱动生成的可控性。

为了兼顾训练效率与实验闭环的可复现性，本文并未采用大规模混合风格数据，而是控制数据规模在适中范围内，并保证样本风格分布尽可能一致。每个样本均由图像和对应文本描述共同构成最小训练单元，从而在数据层面建立“视觉内容—文本语义—风格特征”之间的直接对应关系。该组织方式的优点在于结构清晰、便于批量处理和人工校验，既能够降低工程复杂度，又有助于把研究重点集中在风格学习效果和生成质量分析上。

\subsection{图像数据预处理与元数据导出}

图像预处理模块的核心任务是把来源不一、尺寸各异的原始样本转换为适合扩散模型训练的统一输入。具体而言，系统首先扫描全部样本并校验可读性，剔除损坏文件；随后将有效图像统一转换为RGB格式，并采用中心裁剪结合缩放的方式生成固定分辨率训练图像。由于本文重点关注场景与背景风格学习，中心裁剪策略能够较好地保留主景区域，同时避免因输入尺寸不一致导致的批处理困难。

在样本划分阶段，系统通过固定随机种子打乱样本顺序，再按照预设比例切分训练集与验证集。这种做法既保证了每轮预处理结果的可重复性，也避免了数据按原始命名顺序聚集可能带来的偏差。与许多仅输出处理后图像的简单脚本不同，本文的预处理模块还同步导出元数据文件，记录处理后图像标识及其对应文本描述，供训练阶段的数据加载器直接使用。

除训练元数据外，系统还额外保存原始样本到处理样本的映射关系以及总体统计信息，如扫描数量、有效数量、训练验证划分结果和caption完整性等。这样的设计有两个明显价值：其一，当某一张生成效果异常时，可以快速回溯到对应原始样本；其二，当需要在论文中说明数据组织过程时，可以直接引用统计结果而不依赖人工估算。由此可见，元数据组织本身就是实验可解释性的一部分。

\subsection{视频帧数据预处理设计}

考虑到本文的视频部分主要用于原型验证而非完整视频底模训练，视频数据的组织目标与图像数据并不完全相同。图像数据更关注风格统计分布的稳定性，而视频数据则更关注片段内部帧序与基本可用性。因此，视频预处理模块首先以片段为单位读取连续帧，并根据帧编号排序，保证时序关系不被打乱；随后再依据最小帧数、最大帧数和采样步长等规则筛除不合格片段，保留能够支撑一致性计算的样本。

在帧级处理上，系统与图像阶段保持一致，同样执行RGB转换、中心裁剪和尺寸统一。这样做能够保证图像阶段和视频阶段在空间尺度上的一致性，减少后续递推式img2img生成时因输入分辨率变化带来的不稳定因素。与此同时，视频元数据以片段为单位记录帧序列、帧数和对应文本描述，使系统既可以将其视为连续序列，也可以在需要时回溯到单帧层面。

\section{基于LoRA的风格学习模型设计}

\subsection{训练底模与条件输入设计}

在模型设计上，本文采用Stable Diffusion系列模型作为风格学习底座。该类模型具有成熟的潜空间扩散架构和文本条件生成能力，能够较好地适应小样本风格微调需求。训练时，首先通过VAE将输入图像编码为潜变量，再利用文本编码器将caption转换为条件语义向量，随后由U-Net在给定时间步和文本条件下完成噪声预测。该结构继承了第二章所述的潜空间扩散模型优势，使得系统可以在较高分辨率下仍保持可接受的显存开销和训练效率。

从条件输入角度看，模型训练以“图像潜变量 + 文本条件 + 时间步编码”为基本输入结构。图像首先通过VAE编码到潜空间，再由文本编码器将caption转换为语义向量，U-Net则在指定扩散时间步上预测噪声残差。文本经过tokenizer编码后被截断或填充到固定长度，以便与图像批量输入一同送入模型。在实现中，系统还保留了caption dropout机制，即在一定概率下将文本置空，以提升模型对弱条件输入的鲁棒性。

\subsection{LoRA注入策略与参数配置}

LoRA的基本思想是在冻结原始权重的前提下，用两个低秩矩阵近似原始参数更新量。若某线性层原始权重记为W，则引入LoRA后的权重可表示为W' = W + BA，其中A和 B为待学习的低秩矩阵。本文在实现中将LoRA主要注入到U-Net注意力模块的关键投影层，使风格相关的特征重组过程能够通过少量附加参数被重新塑形。由于注意力层既参与局部纹理组织，也参与文本条件对齐，因此这一注入位置在风格学习任务中具有较高效率。

关键参数包括LoRA秩、学习率、训练步数、梯度累积步数、分辨率、混合精度模式和检查点保存间隔等。其中，LoRA秩决定适配器容量，秩越高，理论上能够表达更复杂的风格偏移，但训练开销也会增加；学习率和训练步数共同决定参数更新强度\cite{ref36}，若设置过小则风格不明显，若设置过大又容易引起底模固有结构被过度改写。本文在多轮实验中采用统一的基础配置，并围绕训练步数和底模先验展开重点比较，这种策略有助于将实验差异收敛到少数关键变量上。

此外，为了兼顾显存和训练时间，系统采用小批量训练配合梯度累积的方式实现较高分辨率训练，并启用fp16混合精度与TF32加速。在工程上，这些设置并不直接改变风格学习原理，却显著提升了实验可执行性，使多组底模比较与长训实验成为可能。

\subsection{训练目标与优化流程}

训练流程延续标准扩散模型的噪声预测框架。设输入图像经VAE编码后得到潜变量z0，在随机时间步t上加入高斯噪声后形成zt，模型则在文本条件c的约束下预测噪声项epsilon。其优化目标可以写为

\begin{equation}
L_{\mathrm{diff}} = \mathbb{E}_{z_0,\epsilon,t}\left[\left\lVert \epsilon - \epsilon_\theta(z_t, t, c) \right\rVert_2^2\right]
\end{equation}

与全参数训练不同，本文在反向传播阶段仅更新LoRA适配器参数，而冻结VAE、文本编码器和U-Net主体权重。这样做使优化过程集中于风格适配而非底模重建，既降低了训练成本，也减少了样本规模有限时发生过拟合崩坏的风险。为了保证训练过程稳定，系统同时引入梯度裁剪、阶段性保存和最佳损失记录等机制，使实验既可中断恢复，也便于比较不同训练阶段对应的生成结果。

需要指出的是，扩散损失的下降仅说明模型在噪声预测意义上更贴近训练样本分布，并不必然意味着主观风格更优。因此，训练优化流程从一开始就与图像对照和可视化报告模块配套设计，即每轮训练结束后都要通过统一推理与对照流程验证效果，而不是仅凭损失值判断训练是否成功。

\section{生成与评估模块实现}

\subsection{单图推理与基线对比实现}

单图推理模块的作用是对训练后权重进行快速验证。系统在加载预训练扩散模型后叠加LoRA权重，并在固定提示词、负向提示词、随机种子和采样参数的条件下输出样例图像。该过程既可用于训练后即时检查，也可用于挑选代表性场景观察风格迁移是否已经出现明显方向。为了避免不同轮次推理结果不可重复，系统会同时保存推理参数摘要，使同一组样例能够被后续精确复现。

更重要的是，本文并未停留在单张样例展示上，而是额外实现了批量提示词对照模块。该模块针对统一的场景提示词集合，先用未加载LoRA的底模生成baseline图像，再在完全相同的种子与参数下加载LoRA生成对应结果。由于随机性已被最大限度控制baseline与LoRA之间的差异就更可以归因于LoRA学到的风格偏移本身。这种批量对照方式显著提升了实验结论的可靠性，也避免了只凭少量“挑图”得出片面判断。

在系统实现上，批量对照结果会被统一整理为比较清单，记录提示词、图像路径和对应实验参数，为后续可视化脚本提供稳定输入接口。由此，图像生成模块不再只是输出图片，而是成为整个结果分析链路的前置环节。

\subsection{短视频生成原型实现}

在视频部分，本文没有直接训练完整视频扩散模型，而是采用递推式短视频原型实现方案。具体而言，系统首先通过文本到图像管线生成首帧，然后把上一帧作为图像到图像管线的输入，在保持提示词主体不变的前提下逐帧生成后续画面，从而构成一段短视频序列。该方案的优点是实现成本低、与已有图像LoRA兼容性强，能够在不额外引入庞大视频主干网络的前提下，初步观察风格能力能否向连续帧迁移\cite{ref37,ref38}。

为了提高递推生成的连续性，系统在视频生成时固定基础随机种子、guidance scale、strength和帧数等关键参数，并分别生成baseline视频和LoRA视频，保证两者可以进行一一对应比较。这样得到的视频序列既能用于定性观察，也能直接送入后续一致性指标计算模块。

\begin{equation}
I_0 = G_{\mathrm{txt2img}}(p, n, s), \quad I_t = G_{\mathrm{img2img}}(I_{t-1}, p, n, s+t, \alpha)
\end{equation}

\subsection{时序一致性评价指标设计}

为了对短视频原型结果进行量化分析，系统实现了无参考时序一致性评估模块。与传统图像评价不同，视频评价更关注相邻帧之间的平滑过渡和结构稳定性，因此系统设计了三类基础指标。第一类是相邻帧像素平均绝对差，用于衡量整体画面变化幅度；第二类是边缘变化平均绝对差\cite{ref39}，通过先计算灰度图梯度，再比较相邻帧的边缘幅值差异，反映结构轮廓的稳定程度；第三类是亮度标准差，用于衡量全局亮度在时间维度上的波动情况。三者共同构成对视频时序一致性的基础描述。

设第t帧图像为I\_t，边缘图为E\_t，帧亮度均值为Y\_t，则本文采用的三个主要指标可表示为

\begin{equation}
\mathrm{MAD} = \frac{1}{T-1}\sum \mathrm{mean}\left(\left|I_t - I_{t-1}\right|\right)
\end{equation}

\begin{equation}
\mathrm{MED} = \frac{1}{T-1}\sum \mathrm{mean}\left(\left|E_t - E_{t-1}\right|\right)
\end{equation}

\begin{equation}
\mathrm{LS} = \mathrm{std}\left(\{\mathrm{mean}(Y_t)\}\right)
\end{equation}

其中，MAD用于描述相邻帧整体变化幅度，MED用于描述边缘结构稳定性，LS用于描述亮度闪烁程度。数值越低，通常意味着帧间过渡越平滑或波动越弱。虽然这组指标不能完全替代更复杂的视频质量评价体系，但对于当前短视频原型验证已足够实用，因为它们能够以统一标准比较baseline与LoRA视频在连续性上的相对变化。

\section{关键算法描述}

\SetKwInput{KwIn}{输入}
\SetKwInput{KwOut}{输出}

\subsection{图像预处理算法}

图像预处理算法的目标是在保持样本主体区域尽量完整的前提下，将原始风格样本变换为统一尺寸的训练输入，并同步导出训练所需的文本元数据与源样本映射。该算法以递归扫描、有效性校验、中心裁剪、固定比例划分和元数据导出为主线，强调输入规范化与结果可回溯。

其算法流程可表述为：

\begin{algorithm}[htbp]
\caption{图像预处理算法}
\KwIn{原始图像目录 $D$，目标分辨率 $R$，验证集比例 $\rho$}
\KwOut{处理后的训练集、验证集及元数据文件}
递归扫描目录 $D$ 中的全部图像文件\;
移除损坏或无法读取的图像文件\;
\ForEach{有效图像 $x_i$}{
  裁剪中心正方形区域，并缩放到 $(R,R)$\;
  若存在同名文本文件，则读取对应 caption\;
}
使用固定随机种子打乱全部有效样本\;
按照比例 $\rho$ 划分得到 $D_{\mathrm{train}}$ 与 $D_{\mathrm{val}}$\;
导出图像文件、metadata jsonl 及源样本映射关系\;
\end{algorithm}

\subsection{LoRA训练算法}

LoRA 风格学习算法以预训练扩散模型为基础，通过冻结主体参数并仅优化低秩适配器，完成特定风格统计分布的高效迁移。该算法的关键在于：一方面保留底模已有的场景构建能力，另一方面利用目标风格数据对注意力投影层进行局部重塑，从而在较低训练成本下获得可观测的风格偏移。

其算法流程可表述为：

\begin{algorithm}[htbp]
\caption{LoRA 风格适配训练算法}
\KwIn{预训练扩散模型 $M$，训练集 $D_{\mathrm{train}}$，LoRA 秩 $r$}
\KwOut{LoRA 权重 $\Delta W$}
冻结 VAE、文本编码器与原始 U-Net 主体参数\;
在注意力投影层中插入 LoRA 适配器\;
\ForEach{训练步骤}{
  将图像 $x$ 编码为潜变量 $z_0$\;
  将 caption $c$ 编码为文本嵌入 $h$\;
  采样时间步 $t$ 与高斯噪声 $\epsilon$\;
  构造带噪潜变量 $z_t$\;
  预测 $\hat{\epsilon}=\epsilon_{\theta}(z_t,t,h)$\;
  计算扩散损失 $L=\lVert \epsilon-\hat{\epsilon}\rVert_2^2$\;
  仅使用 AdamW 更新 LoRA 参数\;
}
保存 LoRA 权重与训练状态\;
\end{algorithm}

\subsection{视频递推生成算法}

递推式短视频生成算法的出发点是在不额外训练视频主干网络的条件下，复用图像风格 LoRA 构造连续帧。其基本策略是先生成首帧，再以前一帧为条件递推生成后一帧，借助图像到图像转换保持场景连续性，并通过统一参数控制使 baseline 与 LoRA 序列具备可比性。

其算法流程可表述为：

\begin{algorithm}[htbp]
\caption{视频递推生成原型算法}
\KwIn{提示词 $p$，负向提示词 $n$，LoRA 适配器 $L$，帧数 $T$}
\KwOut{帧序列 $\{I_t\}_{t=0}^{T-1}$}
加载 text-to-image 与 image-to-image 推理管线\;
按需将 LoRA 适配器 $L$ 注入两条推理管线\;
通过 text-to-image 管线生成首帧 $I_0$\;
\For{$t\leftarrow1$ \KwTo $T-1$}{
  将上一帧 $I_{t-1}$ 作为输入图像\;
  在强度参数 $\alpha$ 下通过 image-to-image 管线生成 $I_t$\;
}
保存全部帧，并组合为 GIF 或视频清单文件\;
\end{algorithm}

\section{本章小结}

本章对本文系统的设计与实现方案进行了说明。首先，围绕研究目标和现实约束，给出了以图像风格学习为核心、以短视频原型验证为扩展的总体方案，并对系统模块进行了清晰划分；其次，介绍了图像样本与视频片段的构建策略和预处理流程，说明数据层如何为后续训练与分析提供稳定基础；再次，阐述了基于Stable Diffusion系列模型的LoRA风格学习方法，包括底模选择、条件输入、注入策略、训练优化和配置驱动实验组织方式；最后，对图像推理、批量对照、可视化绘图、短视频生成原型以及时序一致性评估模块进行了实现层面的梳理，并通过关键算法描述概括了系统运行逻辑。上述设计共同构成了第四章实验与结果分析的实现前提，也说明本文的研究工作已经从方法设想落实为可复现的实验系统。
